\documentclass[11pt,a4paper]{article}
\usepackage[utf8]{inputenc}
\usepackage[T1]{fontenc}
\usepackage{geometry}
\usepackage{booktabs}
\usepackage{tabularx}
\usepackage{enumitem}
\usepackage{hyperref}

\geometry{margin=1in}

\title{\textbf{Definition and Classification of Ecosystem Services}}
\author{}
\date{}

\begin{document}

\maketitle

\section{Core Definition}
\textbf{Ecosystem services} are the direct and indirect benefits that humans obtain from functioning ecosystems and biodiversity. Formalized globally by the \textit{Millennium Ecosystem Assessment} (2005) and expanded by frameworks such as \textit{IPBES} (Nature's Contributions to People), the concept links ecological processes and biophysical structures directly to human well-being, economic security, and survival.

\section{The Four Primary Categories}
Ecosystem services are standardly classified into four functional categories:

\begin{table}[h!]
\centering
\renewcommand{\arraystretch}{1.35}
\begin{tabularx}{\textwidth}{>{\bfseries}l X X}
\toprule
Category & Core Function & Key Examples \\
\midrule
Provisioning & Tangible material or energy outputs extracted directly from ecosystems. & Food crops, timber, freshwater, medicinal plants, genetic resources, biofuels. \\
Regulating & Benefits obtained from the regulation of ecosystem processes and environmental conditions. & Carbon sequestration, climate regulation, flood attenuation, water purification, pollination, erosion control. \\
Supporting & Fundamental ecological processes required to sustain all other ecosystem services (long-term, indirect). & Soil formation, nutrient cycling (N, P, C), primary production (photosynthesis), habitat provision. \\
Cultural & Non-material, intangible benefits that enrich cognitive, spiritual, and social well-being. & Recreation and ecotourism, aesthetic inspiration, cultural heritage, spiritual significance, scientific education. \\
\bottomrule
\end{tabularx}
\caption{Standard classification of ecosystem services based on the Millennium Ecosystem Assessment framework.}
\end{table}

\section{From Ecological Function to Human Value}
In environmental accounting and remote sensing (e.g., mapping via NDVI, biomass, or land-cover metrics), ecosystem services are often analyzed along a cascade:
\begin{enumerate}[label=\textbf{\arabic*.}]
    \item \textbf{Biophysical Structure \& Process:} Underlying ecological components (e.g., forest canopy density, soil organic matter).
    \item \textbf{Ecological Function:} The capacity of the system to perform work (e.g., primary productivity, rainfall interception).
    \item \textbf{Ecosystem Service:} The human-relevant flow derived from that function (e.g., reduced downstream flood risk or harvestable timber).
    \item \textbf{Human Benefit \& Value:} Socioeconomic, health, or intrinsic value realized by society.
\end{enumerate}

\end{document}
